\section{Trabajo Evaluativo}

\img{./images/05/evaluacion}{Evaluación}

\begin{tcolorbox}[]
Objetivos
\tcblower
El siguiente trabajo tiene como objetivo integrar los conocimientos referidos a \emph{Gestión del Mantenimiento} vistos en el transcurso de la primera parte de Mantenimiento y Reparación de Equipos. Para ellos se realizarán actividades que involucren todos los temas vistos, entre ellos: Obsolescencia programada, los tipos y modelos de mantenimiento,codificación de equipos, fichas de equipo, listas de equipos, diagramas de flujo para mantenimiento condicional, asi como manejo de datos estadísticos para modelo de mantenimiento sistemático, analisis de criticidad y los mantenimientos programados.
\end{tcolorbox}

\subsection{Contenido}
El trabajo deberá contener los siguientes puntos listados:
\begin{enumerate}
    \item Investigación (texto mínimo 150 palabras c/u)
    \item Realizar para cada equipo su correspondiente \emph{ficha de equipo}. Que deberá contener:
    \item Realizar lista de equipos donde estén todos los equipos pertenecientes a la planta industrial( mínimo 6).
    \item Anexos y fuentes: Se puede agregar cualquier contenido que sea importante (opcional) y además se debe agregar en la última hoja las fuentes de inforación (obligatorio) siguiendo el formato APA.
\end{enumerate}


\newpage